\documentclass[11pt]{article}
\usepackage{amsmath,amssymb}
\title{Empirical joint-failure counts with uncertain episode links}
\begin{document}
\maketitle

An empirical failure cap can be decidable without deciding which
records belong to each episode. The check requires the possible joins
on one common clock history. A graph made by pooling pairwise matches
from conflicting clock histories can give the wrong answer.

\paragraph{Input and scope.}
Two complete streams each contain \(n>0\) records, one per true episode,
and local binary failure flags \(a_i,b_j\). There are no duplicates,
dropped episodes, or hidden selection in this sample. The first source
provides calibrated acquisition times \(s_i\). The second stores local
acquisition stamps \(t_j\). The model says a join \(\pi\) is admissible
exactly when
\[
 \exists\,\theta\in[\ell,h]\quad \forall i:\quad
          |t_{\pi(i)}-s_i-\theta|\le\epsilon .
\]
There is one common offset \(\theta\) and an independently admissible
bounded residual for each record; no additional chronological or queue
constraints apply. Record identifiers and recorder arrival order are
not episode identities. These are acquisition premises, not guarantees
provided by a recorder or a timestamp format. If the source establishes
a shared drift, batching constraint or order, it must enter the
admissible joins; the model used here does not assume it away for data
that have those constraints.

\paragraph{Exact answer for a cap.}
For a fixed \(\theta\), the admissible pairs form a bipartite graph
\(G_\theta\). Under the stated residual law, its perfect matchings
are exactly the possible joins: a matching and its offset give
the required residual separately for each record. Give edge
\((i,j)\) the cost \(a_i\lor b_j\). Assignment with minimum or maximum
total cost yields the corresponding empirical union count. For the
whole clock interval, cut at every in-range endpoint
\(t_j-s_i\pm\epsilon\). Check every singleton cut and one representative
per open cell, since edge membership is constant within each cell;
in particular, exactly touching intervals can give a feasible
singleton. Take the least and greatest values across the feasible
cells, keeping an actual matching and offset for each bound.
If there are none, the records contradict the acquisition model.
Otherwise a cap on the count is guaranteed iff the upper bound
meets it, and impossible iff the lower bound exceeds it. Multiply
an empirical fraction cap by \(n\) before this test. An empirical
count does not by itself estimate a population probability.

The check is polynomial: there are \(O(n^2)\) cells and assignment
is \(O(n^3)\) per cell in the provided implementation. The replay
audits primal and dual costs. Ordinary weighted assignment is
the direct comparison, not an architecture beaten by this result.
The method avoids choosing a single approximate timestamp join
when more than one is consistent with acquisition.

\paragraph{When can the matching be forgotten?}
Fix a feasible cell and a reference matching \(m\). Contract its
matched pairs to vertices. A feasible \((i,m(j))\) gives directed
arc \(i\to j\), of cost
\[
                    c_{ij}=(a_i\lor b_{m(j)})-(a_i\lor b_{m(i)}).
\]
An alternative matching decomposes into vertex-disjoint alternating
cycles relative to \(m\), and any single directed cycle forms another
matching by replacing its own reference edges. All joins therefore
have the same union count iff every directed cycle has zero total
cost. Equivalently, every strongly connected component admits a
potential \(p\) with \(c_{ij}=p_j-p_i\) on its internal arcs. For
necessity, assign potentials as path sums from a root. Two paths
to a vertex give the same sum by closing them with a return path;
closed walks decompose into cycles. Sufficiency is telescoping.
Arcs between components cannot occur in a replacement cycle.
Within a cell this check is linear in graph size once a reference
matching is supplied, using component labels and propagated
potentials. Across cells the common counts must also agree.
No comparable statement holds for arbitrary future consumers:
on an arbitrary justified graph, the only possible counts can
be 2 and 4, with 3 ruled out. Bounds are exact for monotone cap
decisions, not a promise that every intermediate value occurs.

\paragraph{Replayed cases.}
The pinned recorder callback from the parent experiment carries
authored acquisition stamps and flags as opaque fixture payload.
It also records synthetic transport send/receive times, which
are not used as acquisition evidence. With stamps \(0,2\) in both
streams, residual bound \(1/2\), and offset in \([-2,2]\), the two
cross edges are each possible, but at incompatible offsets.
Their fictitious joint assignment raises the apparent count from
1 to 2. Another case permits multiple joins but always counts
2; changing only the left flags makes 1 and 2 both possible.
A third case has a graph-union matching although there is no
common feasible offset. An endpoint-only control requires exactly
\(\theta=1/2\).

The replay checks 15,924 feasible labeled graph/outcome cases,
including 4,240 with nonunique matching and unique count, and 800
seeded acquisition-model comparisons against direct enumeration
that intersects offset intervals for entire matchings. The
acquisition cases are synthetic; the callback and its publication
fields do not measure clock quality. No physical experiments,
representativeness, or full R4 acceptance follow.
\end{document}
